\medskip

\emph{Cet exercice est un questionnaire à choix multiples.\\ Pour chacune des questions suivantes, une seule des quatre réponses proposées est exacte.\\ Une réponse fausse, une réponse multiple ou l'absence de réponse à une question ne rapporte ni n'enlève de point \\ Pour répondre, indiquer sur la copie le numéro de la question traitée et la lettre de la réponse choisie.\\ Aucune justification n'est demandée.}

\medskip

L'espace est muni d'un repère orthonormé \Oijk{} dans lequel on considère :

\begin{itemize}
\item[$\bullet$~~] les points A$(6~;~-6~;~6)$,  B$(-6~;~0~;~6)$ et C$(-2~;~-2~;~11)$.
\item[$\bullet$~~] la droite $(d)$ orthogonale aux deux droites sécantes (AB) et (BC) et passant par le point A ;
\item[$\bullet$~~] la droite $\left(d'\right)$ de représentation paramétrique :

\[\left\{\begin{array}{l c l}
x &=& -6 - 8t\\
y &=& \phantom{-6 -}4t,\\
z &=& \phantom{-}6 + 5t
\end{array}\right. \text{avec }\:t \in \R.\]
\end{itemize}

\medskip

\textbf{Question 1}

Parmi les vecteurs suivants, lequel est un vecteur directeur de la droite $(d)$ ?

\begin{center}
\begin{tabularx}{\linewidth}{*{4}{X}}
\textbf{a.~} $\vect{u_1}\begin{pmatrix}-6\\3\\0\end{pmatrix}$&
\textbf{b.~} $\vect{u_2}\begin{pmatrix}1\\2\\6\end{pmatrix}$&
\textbf{c.~} $\vect{u_3}\begin{pmatrix}1\\2\\0,2\end{pmatrix}$&
\textbf{d.~} $\vect{u_4}\begin{pmatrix}1\\2\\0\end{pmatrix}$
\end{tabularx}
\end{center}

\medskip

\textbf{Question 2}

Parmi les équations suivantes, laquelle est une représentation paramétrique de la droite (AB) ?

\begin{center}
\begin{tabularx}{\linewidth}{*{2}{X}}
\textbf{a.~}$\left\{\begin{array}{l c r}
x &=& 2t + 6\\
y &=&-6\\
z&=&t+6
\end{array}\right. \:\text{avec }\: t \in \R$&
\textbf{b.~}$\left\{\begin{array}{l c r}
x &=& 2t - 6\\
y &=&- 6\\
z&=&- t - 6
\end{array}\right. \:\text{avec }\: t \in \R$\\
\textbf{c.~}$\left\{\begin{array}{l c r}
x &=& 2t + 6\\
y &=&-t -6\\
z&=&6
\end{array}\right. \:\text{avec }\: t \in \R$&
\textbf{d.~}$\left\{\begin{array}{l c r}
x &=& 2t + 6\\
y &=&t - 6\\
z&=&6
\end{array}\right. \:\text{avec }\: t \in \R$
\end{tabularx}
\end{center}

\medskip

\textbf{Question 3}

Un vecteur directeur de la droite $\left(d'\right)$ est :

\begin{center}
\begin{tabularx}{\linewidth}{*{4}{X}}
\textbf{a.~} $\vect{v_1}\begin{pmatrix}-6\\0\\6\end{pmatrix}$&
\textbf{b.~} $\vect{v_2}\begin{pmatrix}-14\\4\\11\end{pmatrix}$&
\textbf{c.~} $\vect{v_3}\begin{pmatrix}8\\-4\\-5\end{pmatrix}$&
\textbf{d.~} $\vect{v_4}\begin{pmatrix}8\\4\\5\end{pmatrix}$
\end{tabularx}
\end{center}

\medskip

\textbf{Question 4}

Lequel des quatre points suivants appartient à la droite $\left(d'\right)$ ?

\begin{center}
\begin{tabularx}{\linewidth}{*{2}{X}}
\textbf{a.~}M$_1(50~;~-28~;~-29)$&\textbf{b.~}M$_2(-14~;~-4~;~1)$\\
\textbf{c.~}M$_3(2~;~-4~;~-1)$& \textbf{d.~}M$_4 (-3~;~0~;~3)$
\end{tabularx}
\end{center}

\medskip

\textbf{Question 5}

Le plan d'équation $x = 1$ a pour vecteur normal:

\begin{center}
\begin{tabularx}{\linewidth}{*{4}{X}}
\textbf{a.~} $\vect{n_1}\begin{pmatrix}1\\0\\0\end{pmatrix}$&
\textbf{b.~} $\vect{n_2}\begin{pmatrix}0\\1\\1\end{pmatrix}$&
\textbf{c.~} $\vect{n_3}\begin{pmatrix}0\\1\\0\end{pmatrix}$&
\textbf{d.~} $\vect{n_4}\begin{pmatrix}1\\0\\1\end{pmatrix}$
\end{tabularx}
\end{center}
